\subsection{AQ-14-NONPROPER: the rational cohomological implication for all smooth quasiprojective parameters}
Retain the algebraically closed base $k$ and smooth projective $X$.
Let $N^i_{\mathrm{qp}}(X)$ be the relation defined in
AQ-14-QUASIPROJECTIVE, using the \emph{abelian} Albanese of a
smooth connected quasiprojective parameter and free degree-zero
zero-cycles. Put
$P^i_{\mathrm{qp}}(X)=\mathrm{CH}^i(X)_{\mathrm{alg}}/
N^i_{\mathrm{qp}}(X)$. For every prime
$\ell\ne\operatorname{char}k$, one has
\[
 N^i_{\mathrm{qp}}(X)\cap\mathrm{CH}^i(X)[\ell^\infty]
       \subseteq\ker\lambda_X^i.                            \tag{AQ.14.2}
\]
Consequently the implication of original 14/038 holds for this
quasiprojective-parameter quotient in every characteristic:
if $A$ is abelian, $Z\in\mathrm{CH}^i(A\times X)$, and
$u_Z:A(k)\to P^i_{\mathrm{qp}}(X)$ is zero integrally, then
\[
 a({}^tZ):H^{2(\dim X-i)+1}(X,\mathbb Q_\ell(\dim X-i))
                       \longrightarrow H^1(A,\mathbb Q_\ell)
\]
is zero. The same implication holds for the smaller relation
using the semiabelian Albanese kernel. No equality
$N^i_{\mathrm{qp}}=N^i_{\mathrm{proj}}$ in arbitrary positive
characteristic is asserted.

\paragraph{Inputs and coefficient conventions.}
Write $H_j^S(T,R)$ for Suslin homology: it is the homology of
the complex generated in degree $n$ by integral closed
subvarieties of $T\times\Delta^n$ finite and surjective over
$\Delta^n$, tensored with $R$. Its degree-zero cycles are
finite formal sums of closed points. Write $G_T$ for the
pointed Serre semiabelian Albanese and
\[
 0\longrightarrow S_T\longrightarrow G_T\longrightarrow B_T
       \longrightarrow0
\]
for its torus and abelian quotient. Universality shows that
$B_T$ is the abelian Albanese of $T$: a homomorphism from a
torus to an abelian variety is zero. The notation $[\ell^\infty]$
means the subgroup annihilated by some power of $\ell$.

Geisser, \emph{Rojtman's theorem for normal schemes},
arXiv:1402.1831v2, Theorem 1.1, p.~2, gives, for normal $T$,
\[
 H_0^S(T,\mathbb Z)[\ell^\infty]
       \xrightarrow{\sim}G_T(k)[\ell^\infty],\qquad
 H_1^S(T,\mathbb Z)\otimes\mathbb Q_\ell/\mathbb Z_\ell=0.
                                                               \tag{AQ.14.3}
\]
The degree-zero convention on the left is automatic for
torsion. This prime-to-characteristic theorem does not
require a smooth compactification. The finite-coefficient
perfect pairing used below is Geisser--Schmidt (2016),
Theorem 1.1, journal p.~92; their Theorem 7.1, p.~118, identifies it
with the Suslin--Voevodsky comparison. We invoke its
prime-to-characteristic case, where tame and ordinary
abelian $\ell$-primary covers agree. Kahn (2014),
\emph{Algebraic tori as Nisnevich sheaves with transfers},
Lemma 4.1(b), journal p.~707 (July 4, 2013 author version,
Lemma 3(b), p.~8), states
$\operatorname{Alb}_{\mathrm{sa}}(G)=G$ for a semiabelian $G$.

The integral generic-base-change input is Sun,
arXiv:1511.01393v1, Definition 2.2, p.~3, Lemma 2.3(iii),
pp.~4--5, and Theorem 2.5(i), p.~7, with
$\Lambda=\mathbb Z_\ell$. In Lemma 2.3(iii)'s proof, one
dense open works for \emph{all} finite coefficient levels,
by the finite set of Jordan--H\"older constituents. The
constant $\mathbb Z_\ell$ system has the single simple
constituent $\mathbb F_\ell$. Constructibility and preservation
of stratifiability are part of the stated setup and
Remark 2.5.1(ii). We use this only over a finite-type
$\mathbb F_p$-base, $p\ne\ell$, which satisfies the finite
cohomological-dimension condition (LO).

\begin{proof}
Characteristic zero follows already from
AQ-14-QUASIPROJECTIVE and AQ-14-FORWARD. Thus assume
$\operatorname{char}k=p>0$ and $\ell\ne p$.
We first give the three facts needed for specialization.

\emph{1. A correspondence with proper target descends to
Suslin zero-homology.}
For $\Gamma\in\mathrm{CH}^i(X\times T)$, the fibre-evaluation
map $Z_0(T)\to\mathrm{CH}^i(X)$ kills every Suslin boundary.
Indeed, a generator of a Suslin one-chain is an integral
$W\subset\mathbb A^1\times T$ finite and surjective over
$\mathbb A^1$. Its normalization $C$ is a smooth curve,
finite over $W$ and over $\mathbb A^1$. Let $\overline C$
be its smooth projective completion. Pull $\Gamma$ back
along the morphism $X\times C\to X\times T$; pullback is
well defined between smooth varieties by the graph/refined
Gysin construction. Extend a representing cycle by closure
to a class on $X\times\overline C$.

The function $t/(t-1)$ on $\overline C$ has divisor equal
to the fibre over zero minus the fibre over one, including
multiplicities, and no contribution over infinity. The
normalization projection formula identifies their images
in $Z_0(T)$ with the two faces of $W$. Applying the extended
correspondence to this principal divisor gives zero in
$\mathrm{CH}^i(X)$. Its evaluations at points of $C$ agree
with the original fibre evaluations, since refined pullback
is local around those fibres. This proves the claim and
produces a homomorphism
\[
 \Gamma_*:H_0^S(T,\mathbb Z)\longrightarrow\mathrm{CH}^i(X).
                                                               \tag{AQ.14.4}
\]
The construction is compatible with extension of the base
field and with automorphisms of that field. No pushforward
along the nonproper projection $X\times T\to X$ was used.

\emph{2. The integral $H^1$ of the semiabelian Albanese.}
The homology coefficient sequence and (AQ.14.3) give a
natural isomorphism
\[
 H_1^S(T,\mathbb Q_\ell/\mathbb Z_\ell)
             \simeq G_T(k)[\ell^\infty].                    \tag{AQ.14.5}
\]
For clarity, the coefficient sequence has left term
$H_1^S(T,\mathbb Z)\otimes\mathbb Q_\ell/\mathbb Z_\ell$
and right term
$\operatorname{Tor}_1^{\mathbb Z}(H_0^S(T,\mathbb Z),
\mathbb Q_\ell/\mathbb Z_\ell)=H_0^S(T,\mathbb Z)[\ell^\infty]$;
the left term is zero by (AQ.14.3).

The prime-to-$p$ finite pairings of Geisser--Schmidt are
natural and compatible with the coefficient homomorphisms.
Dualizing the direct system
$H_1^S(T,\mathbb Z/\ell^n)$, whose transition is induced by
$\mathbb Z/\ell^n\hookrightarrow\mathbb Z/\ell^{n+1}$,
identifies its dual with the inverse system
$H^1_{\mathrm{et}}(T,\mathbb Z/\ell^n)$ under reduction.
The pairings are viewed in $\mathbb Q_\ell/\mathbb Z_\ell$
via $1\bmod\ell^n\mapsto\ell^{-n}\bmod\mathbb Z_\ell$;
this makes the transition compatibility explicit. Thus
\[
 H^1_{\mathrm{et}}(T,\mathbb Z_\ell)
  \simeq\operatorname{Hom}\bigl(
       G_T(k)[\ell^\infty],\mathbb Q_\ell/\mathbb Z_\ell\bigr)
  \simeq\operatorname{Hom}_{\mathbb Z_\ell}
       (T_\ell G_T,\mathbb Z_\ell).                         \tag{AQ.14.6}
\]
The left side is the inverse limit of the finite groups;
there is no $\varprojlim^1 H^0$ obstruction, because the
constant $H^0$ transition maps are surjective. All
isomorphisms are functorial. Since $G_T$ is its own Serre
Albanese, the Albanese morphism $a_T:T\to G_T$ induces an
isomorphism on $H^1(-,\mathbb Z_\ell)$. This conclusion
is deliberately integral-adic, not an assertion about
$H^1(-,\mathbb Z/\ell^n)$ for each $n$.

\emph{3. The finite-field vanishing.}
Suppose for the moment that $k=\overline{\mathbb F}_p$.
Every class in $H_0^S(T,\mathbb Z)^0$ is torsion.
To see this, join the finitely many points in a representative
by finitely many integral curves in the connected smooth
quasiprojective variety $T$, and pass to their normalizations.
Such curves can be chosen through any prescribed finite
set by hyperplane sections of sufficiently high degree;
only the resulting curves and their normalizations are
used. Differences of points lift to degree-zero classes
on these smooth open curves. For a smooth curve $C$,
$H_0^S(C,\mathbb Z)^0$ is the group of points of its
Rosenlicht generalized Jacobian, equivalently the curve
case of Geisser, Theorem 1.3, p.~3. Each such Jacobian and
each of its points are defined over some finite extension
of a finite field, and its rational points over that
finite field form a finite group. Each lifted class,
and hence each original degree-zero class, is torsion.

Let $v\in H_0^S(T,\mathbb Z)^0$ have zero image in $B_T(k)$,
and let $v_\ell$ denote its $\ell$-primary component.
Under (AQ.14.3), the class $v_\ell$ corresponds to an
element of $S_T(k)[\ell^\infty]$. Consider the homomorphism
\[
 S_T(k)[\ell^\infty]\longrightarrow
 H_0^S(T,\mathbb Z)[\ell^\infty]
 \xrightarrow{\Gamma_*}\mathrm{CH}^i(X)[\ell^\infty]
 \xrightarrow{\lambda_X^i}
 H^{2i-1}(X,\mathbb Q_\ell/\mathbb Z_\ell(i)).                \tag{AQ.14.7}
\]
All varieties, maps and $\Gamma$ descend to a common finite
field $\mathbb F_q$, after enlarging it if needed. Thus
(AQ.14.7) is Frobenius-equivariant. On
$V_\ell S_T=T_\ell S_T\otimes\mathbb Q_\ell$ geometric
Frobenius has eigenvalues $q^{-1}$ times roots of unity:
a torus splits over a finite extension of $\mathbb F_q$.
On
$H^{2i-1}(X,\mathbb Q_\ell(i))$ its eigenvalues have
complex absolute value $q^{-1/2}$, by Deligne (1974),
Theorem (1.6), printed p.~276, after the indicated Tate
twist. These spectra are disjoint. A linear map
commuting with Frobenius between these two spaces is zero:
apply relatively prime characteristic polynomials and
B\'ezout's identity, or use generalized eigenspaces
after extending scalars. Semisimplicity is not
required for this step.

The rational Tate module of the last group in (AQ.14.7)
is exactly $H^{2i-1}(X,\mathbb Q_\ell(i))$. More explicitly,
that group has maximal divisible subgroup
$H^{2i-1}(X,\mathbb Z_\ell(i))_{\mathrm{free}}
 \otimes\mathbb Q_\ell/\mathbb Z_\ell$ and finite quotient
contained in $H^{2i}(X,\mathbb Z_\ell(i))[\ell^\infty]$.
The domain of (AQ.14.7) is divisible. Its image is therefore
in the maximal divisible subgroup, and a homomorphism
between these finite-corank divisible $\ell$-groups is
zero if its map of rational Tate modules is zero.
Consequently (AQ.14.7) is zero, and
$\lambda_X^i(\Gamma_*v_\ell)=0$.

It follows that over $\overline{\mathbb F}_p$ the conclusion
(AQ.14.2) holds: in a finite sum of generators each
$H_0^S$ class is torsion; its prime-to-$\ell$ component
maps to prime-to-$\ell$ torsion, whereas the $\ell$-primary
component has zero Bloch image by the preceding argument.
Taking the $\ell$-primary component of an expression for
$x\in N^i_{\mathrm{qp}}(X)\cap\mathrm{CH}^i(X)[\ell^\infty]$
gives the claimed vanishing.

\emph{4. Spreading to the finite-field situation.}
Now take arbitrary algebraically closed $k$ of characteristic
$p>0$ and a fixed
$x\in N^i_{\mathrm{qp}}(X)\cap\mathrm{CH}^i(X)[\ell^\infty]$.
Choose a finite expression
$x=\sum_j\Gamma_{j*}v_j$, with $v_j\in Z_0(T_j)^0$
mapping to zero in $B_{T_j}(k)$, and choose $r$ with
$\ell^r x=0$. Spread all varieties, finite zero-cycles,
correspondences, rational-equivalence witnesses for these
equalities, and the pointed maps
$T_j\to G_{T_j}\to B_{T_j}$ over an integral finite-type
$\mathbb F_p$-scheme $S$. After shrinking, $S$ is regular,
$X/S$ is smooth projective, the $T_j/S$ are smooth with
geometrically connected fibres, and the $G_j/S$ are
semiabelian schemes in exact sequences
$0\to S_j\to G_j\to B_j\to0$ with tori $S_j$ and
abelian schemes $B_j$. The sections representing the
zero-cycles remain inside $T_j$, and the relations
$B_j(v_j)=0$ persist. The spread maps need not be
relative universal Albanese morphisms.

Apply Sun's integral generic base change simultaneously to
the finitely many $T_j\to S$ and $G_j\to S$. Their
$R^1(-)_*\mathbb Z_\ell$ are constructible adic sheaves,
and become lisse on a further dense open. By (AQ.14.6)
the map induced by $T_j\to G_j$ is an isomorphism at
a geometric generic point. Indeed its base extension
to $k$ is the chosen Albanese map, and prime-to-$p$
\'etale cohomology is unchanged by extension of
algebraically closed fields. Its kernel and cokernel
vanish on a further dense open, so this is an
isomorphism of integral adic sheaves there. Generic
base change now identifies its stalk at \emph{every}
geometric point of that open with the actual map
on $H^1(-,\mathbb Z_\ell)$ of the fibre. This uses one
open for the full integral system, not separate opens
for infinitely many coefficient levels.

Choose a closed point $s$ in this open; its residue
field is finite. Over $\overline{k(s)}$, the map
$T_{j,s}\to G_{j,s}$ factors by universality as
\[
 T_{j,s}\longrightarrow G_{T_{j,s}}
                  \xrightarrow{h_j}G_{j,s}.
\]
By (AQ.14.6), $T_\ell h_j$ is an isomorphism. A
homomorphism of semiabelian varieties whose Tate module
is an isomorphism is an isogeny with no $\ell$-primary
kernel: a positive-dimensional connected kernel or
quotient has a nonzero Tate module, and a nonzero finite
$\ell$-primary kernel makes the Tate-module image a
proper sublattice. The induced homomorphism of abelian
quotients
\[
 b_j:B_{T_{j,s}}\longrightarrow B_{j,s}
\]
is also an isogeny with no $\ell$-primary kernel.
Indeed an isogeny of semiabelian varieties maps the
torus onto the torus; the kernel of its map of abelian
quotients is a quotient of its finite kernel. The first
assertion follows since the resulting quotient of the
target by the image torus is a quotient of the proper
source abelian variety, and hence has zero toric part.

The equality $B_{j,s}(v_{j,s})=0$ and injectivity of
$b_j$ on $\ell$-primary torsion imply that the
$\ell$-primary component of the actual abelian Albanese
class of $v_{j,s}$ is zero. Although the whole class
need not vanish, this is precisely what step 3 needs:
the $\ell$-primary Suslin class of $v_{j,s}$ corresponds
to the actual torus $S_{T_{j,s}}$. Thus the argument
in step 3 gives $\lambda_{X_s}^i(x_s)=0$.

Finally specialize from a geometric generic point to
$s$ along a regular flag through $s$, as in the complete
flag construction in AQ-14-FORWARD. The finite cycle
and rational-equivalence data make $x_s$ the
specialization of $x$, preserve $\ell^r x_s=0$, and
identify each specialized fibre evaluation with
$\Gamma_{j,s*}v_{j,s}$. Refined Gysin functoriality
provides this compatibility locally at the sections.
Bloch's specialization theorem sends
$\lambda_X^i(x)$ to $\lambda_{X_s}^i(x_s)$, and smooth
proper base change for $X/S$ makes the cohomological
specialization an isomorphism. Cohomological invariance
under algebraically closed extension also identifies
the original $k$-fibre with the geometric generic
fibre for this calculation. Therefore
$\lambda_X^i(x)=0$, proving (AQ.14.2).

For the final cohomological implication, lift each
$\ell$-power torsion point of $A$ to its torsion
zero-cycle class by prime-to-$p$ Roitman. The integral
premise puts its image under $Z$ in the left side of
(AQ.14.2). Bloch correspondence compatibility, followed
by Tate modules and tensoring with $\mathbb Q_\ell$,
therefore kills
$Z_*:H^{2\dim A-1}(A,\mathbb Q_\ell(\dim A))
\to H^{2i-1}(X,\mathbb Q_\ell(i))$.
Poincar\'e duality gives exactly the displayed transpose.
This is the same final argument as AQ-14-FORWARD, now
with the broader relation. For the semiabelian-kernel
convention, its relation is contained in the
abelian-kernel relation, so the torsion vanishing and
the implication follow as well.
\end{proof}
